
\usepackage[english,russian]{babel}
\usepackage[T2A]{fontenc}


% Абзацы и геометрия страницы
\usepackage{indentfirst}
\usepackage{geometry}

% Математика
\usepackage{amsmath}
\usepackage{amsthm}
\usepackage{amsfonts}
\usepackage{amssymb}

% Графика
\usepackage{tikz}
\theoremstyle{definition}
\newtheorem{definition}{Определение}
\newtheorem{example}{Пример}
\newtheorem{remark}{Замечание}

\theoremstyle{plain}
\newtheorem{theorem}{Теорема}

\newcommand{\N}{\mathbb{N}}
\newcommand{\R}{\mathbb{R}}
\newcommand{\eps}{\varepsilon}
\newtheorem{axiom}{Аксиома}

\usepackage{cancel}

% Списки и таблицы
\usepackage{enumitem}
\usepackage{tabularx}

% Дополнительные пакеты
\usepackage{blindtext}
\usepackage{titlesec}

% Ссылки
\usepackage[
    colorlinks=true,
    linkcolor=blue,
    urlcolor=blue,
    citecolor=blue
]{hyperref}

% Колонтитулы
\usepackage{fancyhdr}
\pagestyle{fancy}
\fancyhf{}

% Берём название из \section
\renewcommand{\sectionmark}[1]{\markright{#1}}

\fancyhead[L]{\rightmark}
\fancyhead[R]{\hyperref[contents]{Содержание}}

\fancyfoot[C]{\thepage}

\renewcommand{\headrulewidth}{0.4pt}


% Своя стрелка
\newcommand{\myarrow}{%
    \begin{tikzpicture}[baseline=-0.5ex, scale=0.35]
        \draw[->, line width=0.8pt]
            (0,0.5)
            .. controls (-0.2,0.2) and (-0.2,-0.2) ..
            (0.2,-0.2)
            .. controls (0.5,-0.2) and (0.8,-0.2) ..
            (1.2,-0.2);
    \end{tikzpicture}%
}
